%% ---------------------------------------------------------------------
\section{Tipo y error: todo está en las bajas frecuencias}
\label{sec:rf-tipo}
%% ---------------------------------------------------------------------

Con retroalimentación unitaria, el error de seguimiento $E(s)=R(s)-Y(s)$ en estado estacionario depende del número de integradores de la malla abierta (el \emph{tipo}) y de una constante estática (\cref{cap:er}). Las dos cosas se leen en el extremo izquierdo del diagrama de magnitud (\cref{fig:rf-tipos}):

\begin{itemize}
  \item \textbf{Tipo 0:} la magnitud arranca plana. El nivel de baja frecuencia es $20\log\kp$.
  \item \textbf{Tipo 1:} arranca con $-20\dBdec$. La asíntota de baja frecuencia, $\kv/\omega$, cruza $0\dB$ en $\omega=\kv$.
  \item \textbf{Tipo 2:} arranca con $-40\dBdec$. La asíntota, $\ka/\omega^2$, cruza $0\dB$ en $\omega=\sqrt{\ka}$.
\end{itemize}

\begin{figure}[ht]
\centering
\includegraphics[width=\linewidth]{cap10/f_tipos}
\caption{El tipo se ve en la pendiente inicial y la constante estática, en un nivel o en un cruce de la asíntota de baja frecuencia.}
\label{fig:rf-tipos}
\end{figure}

\begin{ejemplo}[label=ej:rf-tipos]{Lea las constantes y el error}
Los tres diagramas de la \cref{fig:rf-tipos} corresponden a $\dfrac{0.01}{s/10+1}$, $\dfrac{1}{s(s/10+1)}$ y $\dfrac{100}{s^2(s/100+1)}$. Determine el tipo, la constante estática y el error ante escalón, rampa y parábola unitarios (los tres lazos cerrados son estables).
\solucion
\begin{center}\small
\begin{tabular}{@{}lccccc@{}}
\toprule
Sistema & lectura en el Bode & constante & escalón & rampa & parábola\\
\midrule
Tipo 0 & nivel $-40\dB$ & $\kp=10^{-40/20}=0.01$ & $\frac{1}{1+\kp}=0.99$ & $\infty$ & $\infty$\\
Tipo 1 & cruce en $\omega=1$ & $\kv=1$ & $0$ & $\frac{1}{\kv}=1$ & $\infty$\\
Tipo 2 & cruce en $\omega=10$ & $\ka=10^2=100$ & $0$ & $0$ & $\frac{1}{\ka}=0.01$\\
\bottomrule
\end{tabular}
\end{center}
\end{ejemplo}

\codigorepo{cap10_respuesta_frecuencial/Clase3_LazoCerrado/C3_Ej2_TipoSistema.m}{tipos}{ejemplo \ref{ej:rf-tipos}}

\begin{ideaclave}
Para mejorar el error hay que \emph{subir la curva en bajas frecuencias}; para el transitorio importa la \emph{zona del cruce}. Un compensador en atraso explota exactamente esta separación (\cref{sec:rf-atraso}).
\end{ideaclave}

\begin{pausa}
¿Por qué en un sistema tipo 2 la asíntota de baja frecuencia cruza $0\dB$ en $\sqrt{\ka}$ y no en $\ka$?
\respuesta{En bajas frecuencias $\Gma(\jw)\approx\ka/(\jw)^2$, cuya magnitud es $\ka/\omega^2$. Igualándola a $1$ resulta $\omega^2=\ka$, es decir, $\omega=\sqrt{\ka}$.}
\end{pausa}
